\documentclass[11pt]{article}
\usepackage[margin=0.95in]{geometry}
\usepackage{amsmath,amssymb}
\usepackage{graphicx}
\usepackage{booktabs,longtable,array}
\usepackage{xcolor}
\usepackage[hidelinks]{hyperref}
\usepackage{enumitem}
\usepackage{caption}
\setlist{itemsep=2pt,topsep=4pt}
\captionsetup{font=small,labelfont=bf}
\setlength{\parskip}{5pt}
\setlength{\parindent}{0pt}
\renewcommand{\arraystretch}{1.15}
\title{Phase 1 Rationale: Testing the Obvious Rescue}
\author{}
\date{trm-delay-robust-scaling $\cdot$ phase rationale $\cdot$ 2026-09-08/09 $\cdot$ 48 runs\\[2pt]Why the phase existed, what inspired its design, and what it contributed. Results are in \texttt{.\allowbreak{}.\allowbreak{}/\allowbreak{}phase\_\allowbreak{}summaries/\allowbreak{}phase1\_\allowbreak{}summary.\allowbreak{}md}; labels are in \texttt{.\allowbreak{}.\allowbreak{}/\allowbreak{}glossary.\allowbreak{}md}.\\[2pt]Research question (the whole project).** Does weight-tied recursion (TRM), used as the actor-critic trunk of a model-free RL agent, benefit continuous control under action delay, and if not, why not?}
\begin{document}
\maketitle
\section*{In one paragraph}



Phase 0 ended with an anecdote, "comparable reward at about 18$\times$ the cost", and the rescue it invites: Pendulum is too easy, and harder tasks will reward the extra computation. Phase 1 existed to test both under control. It turned the anecdote into a measurement (TRM lost every Pendulum seed pair) and put the rescue, H1, to a falsifiable threshold test on a long-horizon, sparse-reward task. The threshold failed. Phase 2's re-analysis later showed the task had not been learned by any trunk, so H1 itself was never really at risk.



\section*{1. Where the project stood going in}



\begin{itemize}
\item A verified delay test bed and a documented TRM port (Phase 0).
\item An uncontrolled observation: with augmentation, MLP, Transformer and TRM reached "comparable reward" on Pendulum, with TRM about 18$\times$ slower to train.
\item H1, pre-registered: TRM's advantage over the MLP and Transformer is non-decreasing in task complexity, and becomes measurable once the horizon is about 5$\times$ Pendulum's and the reward is sparse.
\end{itemize}



\section*{2. Why this phase was made}



Two reasons, one defensive and one exploratory.



\begin{enumerate}
\item \textbf{The anecdote needed a control.} "Comparable" came from a pilot with no seeds paired and no statistic. If TRM was genuinely at parity on Pendulum, a harder task could tip it into an advantage. If it was already behind, the framing had to change.
\item \textbf{The rescue needed a test that could fail.} "Harder tasks will need the extra computation" is unfalsifiable unless "harder" and "measurable" are fixed in advance. H1 fixed both: horizon about 5$\times$ Pendulum's, sparse reward, significance at $\geq$ 5/6 pooled seed pairs.
\end{enumerate}



\section*{3. Design decisions and what inspired them}



\begin{center}\small\begin{tabular}{>{\raggedright\arraybackslash}p{0.276\linewidth}>{\raggedright\arraybackslash}p{0.196\linewidth}>{\raggedright\arraybackslash}p{0.434\linewidth}}
\toprule
Decision & Inspiration & Why it mattered \\
\midrule
MountainCarContinuous-v0 as the harder rung & H1's own threshold: 999-step horizon, one sparse terminal reward & It is the standard continuous-control task with both properties, so it tests the claim as written, not a convenient variant. \\
Four trunks: MLP, Transformer, TRM, LSTM & The trunks already in the codebase, plus an LSTM added in Phase 0 planning & The Transformer is a large non-recursive comparator; the LSTM is the classic recurrent baseline for partial observability. TRM's result could then be read against both "big" and "recurrent". \\
Two conditions: no delay (C0) and fixed d = 1 + queue (C1) & Phase 0's verified augmentation & Shows whether the delay changes the ranking, and doubles the pairs available for the sign test. \\
3 seeds $\times$ 2 conditions pooled to n = 6 & Phase 0's power check & n = 6 is the smallest design in which a clean sweep is significant (p = 0.016). n = 3 alone cannot reach p < 0.125. \\
DDPG held fixed & The codebase's default algorithm & Changing only the trunk isolates the architecture. TD3 waited until Phase 4. \\
Wall-clock logged alongside reward & Cost as a first-class metric (project brief) & Turned "18$\times$ slower" into a measured multiple per environment (6.7$\times$ to 12.8$\times$). \\
Re-submission when logging was found lacking (v1 $\rightarrow$ v3) & Operational lessons during the phase & v2 added per-episode logging and v3 step histories. That telemetry is what made Phase 2's re-analysis possible without new compute. \\
\bottomrule
\end{tabular}\end{center}



\section*{4. How it contributed to the research question}



\begin{center}\small\begin{tabular}{>{\raggedright\arraybackslash}p{0.319\linewidth}>{\raggedright\arraybackslash}p{0.615\linewidth}}
\toprule
Before Phase 1 & After Phase 1 \\
\midrule
"Comparable reward" on Pendulum (anecdote) & A consistent deficit: TRM $-$ MLP 0/6 on Pendulum, p = 0.016; $-$37.1 under fixed delay, about 3\% of the span \\
H1's threshold untested & Threshold falsified: TRM $-$ MLP 3/6 on MountainCar, p = 0.66 \\
Cost "about 18$\times$" & 6.7$\times$ to 12.8$\times$ the MLP's wall-clock, growing with task complexity \\
\bottomrule
\end{tabular}\end{center}



For the "does it benefit" half of the question, Phase 1 gave the first controlled answer: no, on the one task every trunk learned. For the "why not" half it opened four competing explanations it could not separate: recursion itself, capacity (TRM had 5.8$\times$ the MLP's parameters), the latent reset, and a task with nothing for recursion to infer.



\section*{5. How it shaped the experiments that followed}



\begin{itemize}
\item \textbf{Phase 2} added a second kind of difficulty (BipedalWalker's coupled dynamics), so a null could not be dismissed as "the wrong kind of hard", plus structure and augmentation ablations requested after Phase 1 review.
\item \textbf{The Pendulum deficit} became the project's anchor contrast. Phase 3 reproduced it at n = 12 ($-$34.4, 0/12, p = 0.0005) and then took it apart.
\item \textbf{The LSTM's failure on Pendulum} ($\approx$ $-$1460) was flagged as a tuning artefact, which later fed Phase 3's decision to use previous-action input and stored-state replay for recurrent trunks (Ni et al., 2022; Kapturowski et al., 2019).
\end{itemize}



\section*{6. In hindsight}



\begin{itemize}
\item \textbf{No budget pilot.} Nobody checked that MountainCar was learnable by the baseline at 200 episodes. It was not: the median run never reached the goal. The narrowing gap was all trunks converging on the same floor. Phase 4 made a budget pilot with a pre-registered "learned" bar mandatory before any complexity matrix.
\item \textbf{The anecdote was overturned cleanly} because the statistic was fixed before the data. That is the main methodological win of the phase.
\end{itemize}



\section*{Sources}



\begin{itemize}
\item Matheron et al. (2019), the DDPG + OU-noise exploration failure (identified post hoc in Phase 2).
\item Ni et al. (2022), recurrent model-free RL as a POMDP baseline.
\item Project files:
\begin{itemize}
\item \texttt{phase0-literature/\allowbreak{}ideation/\allowbreak{}proposal-final.\allowbreak{}json} (H1)
\item \texttt{.\allowbreak{}.\allowbreak{}/\allowbreak{}full\_\allowbreak{}report.\allowbreak{}md} \S{}5
\item \texttt{experiments/\allowbreak{}analysis\_\allowbreak{}v3\_\allowbreak{}final.\allowbreak{}json}
\end{itemize}
\end{itemize}
\end{document}
